\documentclass[11pt,a4paper]{article}
\usepackage[brazil]{babel}
\usepackage[utf8]{inputenc}
\usepackage[T1]{fontenc}
\usepackage{lmodern}
\renewcommand{\baselinestretch}{1.05}
\usepackage{mhchem}
\usepackage{graphicx}
\usepackage{svg}
\usepackage{amssymb}
\usepackage{amsmath}
\usepackage{enumitem}
\usepackage{geometry}
\usepackage{booktabs}
\geometry{top=2cm, bottom=2cm, left=2.5cm, right=2.5cm}

\newcommand{\oC}{$^\circ$C}
\newcommand{\e}[1]{$\times 10^{#1}$}

\setlist[itemize,1]{leftmargin=*, itemsep=2pt, topsep=2pt}
\setlist[enumerate,1]{leftmargin=*, itemsep=2pt, topsep=2pt}

\pagestyle{empty}

\begin{document}

\begin{center}
  {\LARGE\bfseries QG101 -- Química Geral}\\[4pt]
  {\large\bfseries Bloco 13 -- Termoquímica Básica}\\[2pt]
  \textit{Material didático -- Aula Teórica 13}
\end{center}

\bigskip

% ================================================================
\section*{1. Energia e transformações químicas}

Toda transformação química envolve troca de energia entre o
\textbf{sistema} (as substâncias que participam da reação) e a
\textbf{vizinhança} (tudo o que está em volta). Essa troca pode ocorrer
na forma de \textbf{calor} ($q$) e/ou \textbf{trabalho} ($w$).

A \textbf{Termoquímica} estuda especificamente o calor liberado ou
absorvido pelas reações químicas, geralmente sob pressão constante (a
pressão atmosférica, na maioria dos experimentos de laboratório e do
cotidiano).

% ================================================================
\section*{2. Entalpia e variação de entalpia}

A \textbf{entalpia} ($H$) é uma propriedade do sistema relacionada à sua
energia interna ($U$) e ao trabalho de expansão contra a pressão externa:
$H = U + PV$. Como $H$ é uma \textbf{função de estado}, o que importa nas
reações químicas é sua \textbf{variação}:
$$\Delta H = H_{\text{produtos}} - H_{\text{reagentes}}$$

Em um processo à pressão constante, $\Delta H$ é exatamente igual ao calor
trocado com a vizinhança: $\Delta H = q_p$. Por isso $\Delta H$ é a
grandeza termoquímica mais usada no dia a dia do laboratório.

% ================================================================
\section*{3. Reações exotérmicas e endotérmicas}

\begin{itemize}
\item \textbf{Exotérmica} ($\Delta H < 0$): o sistema \textbf{libera}
  calor para a vizinhança, que se aquece. A entalpia dos produtos é
  \textbf{menor} que a dos reagentes. Exemplos: combustões, a maioria das
  reações de neutralização, respiração celular.
\item \textbf{Endotérmica} ($\Delta H > 0$): o sistema \textbf{absorve}
  calor da vizinhança, que se esfria. A entalpia dos produtos é
  \textbf{maior} que a dos reagentes. Exemplos: fotossíntese, fusão do
  gelo, muitas dissoluções salinas (ex.: \ce{NH4NO3} em água).
\end{itemize}

\begin{center}
\includesvg[width=0.92\textwidth]{figuras/bloco13_fig1_diagrama_energia}

{\small Figura 1 -- Diagramas de energia para reações exotérmica e
endotérmica.}
\end{center}

% ================================================================
\section*{4. Calor específico e capacidade calorífica}

O \textbf{calor específico} ($c$) é a quantidade de calor necessária para
elevar em $1$\oC{} a temperatura de $1$~g de uma substância. A relação
entre calor trocado, massa e variação de temperatura é:
$$q = m \, c \, \Delta T$$

A \textbf{capacidade calorífica} ($C = m \, c$) refere-se a uma amostra
específica (não à substância em geral) e tem unidade J/\oC.

A água tem calor específico excepcionalmente alto
($c = 4{,}18$~J~g$^{-1}$\oC$^{-1}$), bem maior que o de metais comuns como
o ferro ($c \approx 0{,}45$~J~g$^{-1}$\oC$^{-1}$). Isso significa que, para
a mesma massa e o mesmo calor fornecido, a água esquenta (ou esfria) muito
mais lentamente -- por isso a areia da praia fica escaldante ao meio-dia
enquanto o mar permanece ameno, e por isso a água é usada em sistemas de
resfriamento e na regulação da temperatura corporal.

% ================================================================
\section*{5. Primeira Lei da Termodinâmica}

A energia não pode ser criada nem destruída, apenas transferida ou
transformada. Para um sistema fechado, isso se expressa como:
$$\Delta U = q + w$$
em que $q$ é o calor \textbf{recebido} pelo sistema (positivo se
absorvido, negativo se liberado) e $w$ é o trabalho realizado
\textbf{sobre} o sistema (positivo se a vizinhança realiza trabalho sobre
o sistema; negativo se o sistema realiza trabalho de expansão sobre a
vizinhança, como ao empurrar o ar atmosférico durante a liberação de um
gás em uma combustão).

\textbf{Exemplo:} se um sistema absorve $300$~J de calor ($q = +300$~J) e
se expande, realizando $120$~J de trabalho sobre a vizinhança
($w = -120$~J), então:
$$\Delta U = 300 + (-120) = +180~\text{J}$$

\clearpage

% ================================================================
\section*{6. Equações termoquímicas}

Uma equação termoquímica é uma equação balanceada acompanhada do valor de
$\Delta H$ correspondente. Três regras práticas:

\begin{enumerate}
\item $\Delta H$ refere-se exatamente às quantidades (em mol) indicadas
  pelos coeficientes estequiométricos da equação.
\item Se a equação for \textbf{invertida}, o sinal de $\Delta H$ também se
  inverte.
\item Se a equação for \textbf{multiplicada} por um fator $k$, $\Delta H$
  é multiplicado pelo mesmo fator $k$.
\end{enumerate}

O \textbf{estado físico} de cada substância também importa: a combustão do
metano formando água \textbf{líquida} libera mais energia do que a mesma
combustão formando água \textbf{vapor}, pois a condensação do vapor de
água libera calor adicional:
$$\ce{CH4(g) + 2O2(g) -> CO2(g) + 2H2O(l)} \qquad \Delta H = -890~\text{kJ/mol}$$
$$\ce{CH4(g) + 2O2(g) -> CO2(g) + 2H2O(g)} \qquad \Delta H = -802~\text{kJ/mol}$$

% ================================================================
\section*{7. Lei de Hess}

Como $H$ é uma função de estado, $\Delta H$ de uma reação \textbf{não
depende do caminho} percorrido -- apenas dos estados inicial e final. A
\textbf{Lei de Hess} afirma que, se uma reação puder ser escrita como a
soma de duas ou mais etapas, seu $\Delta H$ é igual à \textbf{soma} dos
$\Delta H$ dessas etapas.

\textbf{Exemplo:} qual é o $\Delta H$ de formação do \ce{CO} a partir de
\ce{C} e \ce{O2}, dado que:
\begin{align*}
\ce{C(s) + O2(g) &-> CO2(g)} &\Delta H_1 &= -394~\text{kJ/mol (caminho direto)}\\
\ce{CO(g) + 1/2 O2(g) &-> CO2(g)} &\Delta H_3 &= -283~\text{kJ/mol}
\end{align*}

A reação procurada, $\ce{C(s) + 1/2 O2(g) -> CO(g)}$ ($\Delta H_2$), somada
à segunda equação dada reproduz a primeira. Logo:
$$\Delta H_1 = \Delta H_2 + \Delta H_3 \quad\Longrightarrow\quad
\Delta H_2 = \Delta H_1 - \Delta H_3 = -394 - (-283) = -111~\text{kJ/mol}$$

\begin{center}
\includesvg[width=0.62\textwidth]{figuras/bloco13_fig2_lei_hess}

{\small Figura 2 -- Ciclo de Lei de Hess: o caminho direto
($\Delta H_1$) é igual à soma dos caminhos indiretos ($\Delta H_2 + \Delta H_3$).}
\end{center}

% ================================================================
\section*{8. Síntese}

\begin{center}
\begin{tabular}{p{3.2cm}p{5.4cm}p{5.6cm}}
\toprule
Conceito & Ideia central & Fórmula / critério \\
\midrule
Entalpia & Energia trocada à pressão constante & $\Delta H = H_{\text{prod}} - H_{\text{reag}}$ \\
Exo/endotérmica & Direção do fluxo de calor & $\Delta H < 0$ / $\Delta H > 0$ \\
Calor específico & Energia para variar $1$\oC{} de $1$~g & $q = mc\Delta T$ \\
1.ª Lei & Conservação de energia & $\Delta U = q + w$ \\
Equação termoquímica & $\Delta H$ depende de coeficientes e estados físicos & inverter $\Rightarrow$ troca sinal; multiplicar $\Rightarrow$ escala \\
Lei de Hess & $\Delta H$ não depende do caminho & $\Delta H_{\text{total}} = \sum \Delta H_{\text{etapas}}$ \\
\bottomrule
\end{tabular}
\end{center}

\bigskip
\noindent\textit{A lista de exercícios correspondente a este bloco está no
arquivo \texttt{bloco13\_lista\_exercicios.pdf}.}

\end{document}
